% =============================================================================
\section{Additional Dataset and Evaluation Notes}\label{app:dataset}
% =============================================================================
% SUPPLEMENT-CANDIDATE: this appendix can be moved intact to supplementary material.
\subsection{Relationship to the Public fastMRI Prostate Release}
The public fastMRI Prostate dataset contains 312 patients and both T2-weighted and diffusion-weighted data~\cite{Tibrewala2024FastMRIProstate}. The present manifest contains 270 patients because it is a prepared experimental subset. All cohort counts in the main paper refer to this prepared manifest rather than to the full public release.

\begin{table*}[t]
\centering
\caption{Sequence-scope rationale for the present study.}
\label{tab:t2-scope}
\scriptsize
\begin{tabularx}{\textwidth}{L{2.4cm}L{5.2cm}X}
\toprule
Aspect & T2-weighted branch & Diffusion branch in the public dataset \\
\midrule
Clinical information & Zonal anatomy, gland structure, lesion morphology, local extension features~\cite{Turkbey2019PIRADSv21} & Diffusion restriction represented through multiple $b$-values, ADC, and high-$b$-value contrasts \\
Acquisition/\linebreak reconstruction path & Axial T2-weighted turbo-spin-echo with multicoil complex data & EPI diffusion acquisition with directions, repeated averages, gridding/reconstruction, and derived maps~\cite{Tibrewala2024FastMRIProstate} \\
Reason for current scope & Provides a single-sequence complex input with clear image-domain anatomy and minimal modality-fusion confounding & Scientifically important but introduces additional sequence-specific transformations that are outside the present controlled comparison \\
Interpretation & Controlled T2 slice-classification experiment & Not evaluated here; omission is a scope limitation rather than evidence against diffusion information \\
\bottomrule
\end{tabularx}
\end{table*}

\subsection{Slice Labels and Patient Correlation}
The source release assigns T2 and DWI labels on a slice-by-slice basis in isolation~\cite{Tibrewala2024FastMRIProstate}. This increases the number of supervised samples but means that adjacent slices from one patient are not independent observations. The present patient-disjoint partitioning prevents direct train/test leakage. Nevertheless, the reported AUC and AP remain slice-level metrics; confidence intervals that ignore within-patient clustering would understate finite-cohort uncertainty.

\subsection{Why Average Precision Is Reported}
The test cohort contains 63 positive slices among 1,245 total slices, a prevalence of approximately 5.06\%. ROC AUC is useful for threshold-independent ranking but can remain high when positive predictive performance is modest in an imbalanced population. AP summarizes the precision-recall tradeoff over score thresholds and is therefore reported alongside AUC. Threshold-dependent sensitivity, specificity, balanced accuracy, precision, and F1 use the threshold selected only from validation data.

\subsection{Scope of the Fourier-Domain Arm}
The Fourier-domain complex input is generated from the already prepared and phase-aligned complex image using a centered orthonormal FFT. It is not a direct reload of the original scanner acquisition. Consequently, this arm isolates the effect of presenting an invertible Fourier coordinate representation to the same CNN family under the current preprocessing; it does not reproduce the full raw-data reconstruction pipeline of the source dataset.
